% ============================================================================
% Chapitre 2 — Conception et architecture
% ============================================================================
\chapter{Conception et architecture}

Ce chapitre présente la conception du système Tamwil~AI. Il décrit d'abord l'architecture globale en trois couches, puis les modules fonctionnels qui composent le \textit{backend}, le modèle de données sous-jacent et le détail du pipeline \ac{RAG}. Il se conclut par le \textit{mapping} entre les sources de données et les fonctionnalités exposées à l'utilisateur.

\section{Architecture globale}

L'architecture de Tamwil~AI suit une conception modulaire en trois couches~: interface utilisateur, service applicatif, couche de données. La figure~\ref{fig:architecture} en donne une représentation synthétique.

% ============================================================================
% figures/architecture.tex
% Architecture globale — diagramme dessiné manuellement (coordonnées en cm).
% Aucune génération automatique : chaque noeud est placé à la main.
% ============================================================================
\begin{figure}[H]
\centering
\begin{tikzpicture}[
  font=\footnotesize\sffamily,
  every node/.style={align=center, rounded corners=2pt, inner sep=5pt},
  layerband/.style={draw=black!15, fill=black!3, rounded corners=6pt},
  ui/.style={draw=blue!40!black, fill=blue!10, thick, minimum width=3.6cm, minimum height=1.1cm},
  back/.style={draw=green!50!black, fill=green!10, thick, minimum width=3.6cm, minimum height=1.1cm},
  data/.style={draw=orange!60!black, fill=orange!15, thick, minimum width=3.6cm, minimum height=1.1cm},
  ext/.style={draw=red!50!black, fill=red!8, thick, dashed, minimum width=2.6cm, minimum height=0.9cm},
  mod/.style={draw=green!50!black, fill=green!5, thin, minimum width=2.0cm, minimum height=0.8cm},
  arrow/.style={-Stealth, thick, gray!70!black},
]

% Bandeaux de couche (dessinés en arrière-plan)
\begin{scope}[on background layer]
  \node[layerband, minimum width=13cm, minimum height=1.8cm] (UIband) at (0, 5.0) {};
  \node[layerband, minimum width=13cm, minimum height=3.6cm] (BEband) at (0, 1.6) {};
  \node[layerband, minimum width=13cm, minimum height=1.8cm] (DATAband) at (0, -2.1) {};
\end{scope}

\node[anchor=west, font=\scriptsize\sffamily\itshape] at ([xshift=0.25cm] UIband.west) {Interface utilisateur};
\node[anchor=west, font=\scriptsize\sffamily\itshape] at ([xshift=0.25cm, yshift=1.4cm] BEband.west) {Service applicatif};
\node[anchor=west, font=\scriptsize\sffamily\itshape] at ([xshift=0.25cm] DATAband.west) {Couche de données};

% Couche UI
\node[ui] (front) at (0, 5.0) {\textbf{Next.js 16 + React 19}\\[-1pt]Chat SSE, sidebar, profil, dark mode};

% Serveur FastAPI
\node[back] (api) at (0, 2.9) {\textbf{Serveur FastAPI}\\[-1pt]\code{/api/chat} \textbullet{} \code{/api/chat/stream}};

% Routeur
\node[back, minimum width=6.8cm] (router) at (0, 1.5) {\textbf{Routeur d'intentions}\\[-1pt]mots-clés + regex \textbullet{} extraction profil};

% 5 modules métier + 1 utilitaire
\node[mod] (m1) at (-5.1, 0.1) {\textbf{Scoring}\\\scriptsize fundability};
\node[mod] (m2) at (-3.05, 0.1) {\textbf{Diagnostic}\\\scriptsize KPI};
\node[mod] (m3) at (-1.0, 0.1) {\textbf{Matching}\\\scriptsize investisseurs};
\node[mod] (m4) at (1.05, 0.1) {\textbf{Matching}\\\scriptsize subventions};
\node[mod] (m5) at (3.1, 0.1) {\textbf{Q\&A RAG}\\\scriptsize streaming};
\node[mod, fill=green!2, dashed] (m6) at (5.15, 0.1) {\textit{Explication}\\\scriptsize métriques (routeur)};

% Couche données
\node[data] (chroma) at (-3.6, -2.1) {\textbf{ChromaDB persistant}\\[-1pt]\code{chroma\_db/} (9.4 MB)};
\node[data] (json) at (0.2, -2.1) {\textbf{Dataset JSON}\\[-1pt]investisseurs, subventions, métriques, FAQ};
\node[data] (md) at (4.1, -2.1) {\textbf{Guides \& régulations}\\[-1pt]10 documents Markdown};

% API externe
\node[ext] (openai) at (5.9, 2.9) {OpenAI API\\\scriptsize GPT-3.5-turbo};

% Flèches
\draw[arrow] (front) -- node[right, pos=0.5, font=\tiny] {HTTP / SSE} (api);
\draw[arrow] (api) -- (router);
\draw[arrow] (router.south west) -- (m1.north);
\draw[arrow] (router.south) -- (m2.north);
\draw[arrow] (router.south) -- (m3.north);
\draw[arrow] (router.south) -- (m4.north);
\draw[arrow] (router.south east) -- (m5.north);
\draw[arrow, dashed] (router.east) to[bend left=30] (m6.north);

\draw[arrow] (m5) -- (chroma);
\draw[arrow] (m1) -- (json);
\draw[arrow] (m2) -- (json);
\draw[arrow] (m3) -- (json);
\draw[arrow] (m4) -- (json);
\draw[arrow, dashed] (m5.south east) -- (md.north);

\draw[arrow, dashed] (m5.east) to[bend left=20] node[above, font=\tiny] {requête LLM} (openai.west);
\draw[arrow, dashed] (openai.west) to[bend left=20] node[below, font=\tiny] {tokens} (m5.north east);

\end{tikzpicture}
\caption{Architecture globale de Tamwil~AI --- trois couches (interface Next.js, service FastAPI avec routeur et modules, données ChromaDB + fichiers). Les flèches en pointillé représentent les flux optionnels (API externe, module porté par le routeur).}
\label{fig:architecture}
\end{figure}


La couche \textbf{interface} est une application web Next.js qui consomme l'\ac{API} du \textit{backend} en REST et en \ac{SSE}. La couche \textbf{service} est un serveur FastAPI qui expose un point d'entrée unique pour chaque requête utilisateur~; un \textit{routeur} d'intentions détermine quel module traite la demande et, le cas échéant, enrichit la requête avec les informations extraites du profil. La couche \textbf{données} regroupe la base vectorielle ChromaDB (persistée sur disque) et les fichiers \ac{JSON}/Markdown de référence.

Seul un composant externe reste sollicité à chaque requête~\ac{RAG}~: l'API d'OpenAI, utilisée pour la génération de la réponse finale. Tout le reste --- embeddings, stockage, règles métier --- s'exécute localement sur la machine hôte.

\section{Modules fonctionnels}

Le \textit{backend} expose \textbf{cinq modules métier} et un utilitaire d'explication des métriques implémenté directement dans le routeur. Cette distinction entre modules (scoring, diagnostic, matching investisseurs, matching subventions, \ac{RAG} Q\&A) et utilitaires (explication métrique) reflète fidèlement l'organisation du dépôt~: seuls les cinq premiers correspondent à un fichier dédié dans \code{backend/app/modules/}~; l'explication de métriques, qui ne fait qu'extraire et reformuler des champs du fichier \code{metrics\_reference.json}, n'a pas été jugée suffisante pour justifier un module isolé.

\paragraph{Module 1 --- Scoring de \textit{fundability}.} Fonctionnalité phare du projet~: elle évalue la préparation d'une startup à la levée de fonds.
\begin{itemize}
  \item \textbf{Entrée}~: profil startup (stade) + \ac{KPI} (\ac{MRR}, \textit{burn rate}, \textit{churn}, \ac{CAC}, \ac{LTV}\ldots).
  \item \textbf{Données utilisées}~: \code{metrics\_reference.json} (\textit{benchmarks} pondérés par stade).
  \item \textbf{Logique}~:
  \begin{enumerate}
    \item comparer chaque métrique au \textit{benchmark} du stade correspondant (\textit{bad}/\textit{ok}/\textit{good}/\textit{excellent})~;
    \item calculer un score pondéré sur~100 en utilisant le champ \code{weight\_in\_scoring}~;
    \item gérer spécifiquement les métriques inversées (\ac{CAC}, \textit{churn}, \textit{burn rate}, pour lesquelles \textit{plus bas} = meilleur)~;
    \item identifier forces et faiblesses puis générer un plan d'action tiré des \code{improvement\_tips}.
  \end{enumerate}
  \item \textbf{Sortie}~: score sur 100, ventilation par métrique, axes d'amélioration et plan d'action.
\end{itemize}

\paragraph{Module 2 --- Diagnostic financier.} Analyse individuelle des \ac{KPI} fournis, puis génération d'un résumé global selon la distribution des niveaux constatés. Chaque métrique est restituée avec sa définition, sa formule, les \textit{benchmarks} du stade, une interprétation et jusqu'à trois conseils d'amélioration.

\paragraph{Module 3 --- Matching investisseurs.} À partir du profil (secteur, stade, géographie, ticket recherché), le module calcule un score de \textit{matching} sur~4 (un point par critère satisfait), trie les investisseurs par score décroissant et retourne le \textit{top-5} accompagné des raisons du \textit{matching}, du portefeuille et des contacts.

\paragraph{Module 4 --- Matching subventions.} Filtrage séquentiel par pays, stade, secteur (avec gestion d'un \textit{wildcard} \og tous secteurs \fg) et âge maximum de l'entreprise. Les subventions éligibles sont retournées avec le détail des conditions et, le cas échéant, la \textit{deadline} ou la mention \textit{rolling}.

\paragraph{Module 5 --- Q\&A conversationnel (\ac{RAG}).} Chargé des questions ouvertes. Il enrichit la requête avec les deux derniers échanges, récupère les top-$k$ chunks depuis ChromaDB, puis construit un \textit{prompt} injecté dans le \ac{LLM} avec génération en \textit{streaming} \ac{SSE}. L'attribution des sources se fait via \code{source\_urls.py}.

\paragraph{Utilitaire --- Explication des métriques.} Lorsqu'un message contient un terme de métrique (\ac{MRR}, \ac{CAC}, \textit{churn}\ldots), le routeur identifie les métriques mentionnées et restitue définition, formule, interprétation et \textit{benchmarks} à partir de \code{metrics\_reference.json}. Cette fonctionnalité n'est pas un module au sens fichier~: elle est portée par le routeur (fonction \code{\_explain\_metrics()}) et partage ses données avec les modules 1 et 2.

\section{Modèle de données}

Le dataset a été constitué manuellement pour garantir la pertinence des données au contexte marocain et francophone. Il est organisé selon deux familles~: des fichiers \ac{JSON} fortement structurés et des documents Markdown de long format. Le tableau~\ref{tab:dataset} en donne la composition.

\begin{table}[H]
\centering
\caption{Composition du dataset de Tamwil~AI.}
\label{tab:dataset}
\begin{tabularx}{\textwidth}{@{}l l r X@{}}
\toprule
\textbf{Fichier / dossier} & \textbf{Format} & \textbf{Éléments} & \textbf{Contenu} \\
\midrule
\code{investors.json}            & JSON & 41 profils     & VC, BA, accélérateurs, fonds publics, CVC \\
\code{grants.json}               & JSON & 25 programmes  & Subventions Maroc, France, International \\
\code{metrics\_reference.json}   & JSON & 18 métriques   & \ac{KPI} avec \textit{benchmarks} globaux et MAD \\
\code{faq.json}                  & JSON & 54 Q\&A        & 10 catégories (levée, métriques, aides, régulation\ldots) \\
\code{fundraising\_guide/*.md}   & MD   & 6 guides       & Types de financement, pitch deck, term sheet\ldots \\
\code{regulations/*.md}          & MD   & 4 documents    & Création d'entreprise, fintech, fiscalité, RGPD \\
\bottomrule
\end{tabularx}
\end{table}

Les fichiers \ac{JSON} servent aux modules déterministes (\textit{scoring}, diagnostic, \textit{matching}) et alimentent également la base vectorielle pour le \ac{RAG}. Les documents Markdown ne sont utilisés que par le \ac{RAG}. La structure complète d'un profil investisseur, d'une subvention, d'une métrique et d'une \ac{FAQ} est donnée en annexe~\ref{ann:dataset}.

\section{Pipeline RAG}

Le pipeline \ac{RAG} de Tamwil~AI reprend exactement la structure à six étapes décrite au chapitre précédent (figure~\ref{fig:rag_pipeline}). Les choix techniques retenus à chaque étape sont les suivants~:

\begin{itemize}
  \item \textbf{Ingestion}~: six \textit{loaders} --- un par type de document --- construisent des objets \code{Document} \textit{LangChain} avec un champ \code{metadata} homogène (\code{source}, \code{type}, \code{categorie}).
  \item \textbf{Chunking}~: \code{RecursiveCharacterTextSplitter} avec séparateurs hiérarchiques (\code{["\textbackslash{}n\#\# ", "\textbackslash{}n\#\#\# ", "\textbackslash{}n\textbackslash{}n", "\textbackslash{}n", ". ", " ", ""]}), \code{chunk\_size~=~500}, \code{overlap~=~50}.
  \item \textbf{Embedding}~: \texttt{paraphrase-multilingual-MiniLM-L12-v2}, exécuté localement, 384 dimensions, support français natif.
  \item \textbf{Indexation}~: ChromaDB en mode persistant, collection \code{tamwil\_knowledge\_base}, \textit{backend} SQLite. Taille mesurée après indexation~: environ 9,4~MB.
  \item \textbf{Retrieval}~: recherche par similarité cosinus, $k = 5$ par défaut, filtrage optionnel par \code{filter\_type} pour restreindre la recherche à un type de document donné.
  \item \textbf{Generation}~: appel au modèle OpenAI \textit{GPT-3.5-turbo} avec température 0,3 et mode \textit{streaming}~; si la clé \ac{API} est absente, un mode \textit{fallback} retourne directement les chunks récupérés.
\end{itemize}

\section{Mapping dataset \texorpdfstring{$\to$}{->} fonctionnalités}

Chaque source alimente une ou plusieurs fonctionnalités, comme indiqué dans le tableau~\ref{tab:mapping}.

\begin{table}[H]
\centering
\caption{Mapping entre sources de données et fonctionnalités exposées.}
\label{tab:mapping}
\begin{tabularx}{\textwidth}{@{}l X X@{}}
\toprule
\textbf{Fonctionnalité} & \textbf{JSON principal} & \textbf{Markdown (RAG)} \\
\midrule
Scoring                & \code{metrics\_reference.json}  & --- \\
Diagnostic             & \code{metrics\_reference.json}  & --- \\
Explication métriques  & \code{metrics\_reference.json}  & --- \\
Matching investisseurs & \code{investors.json}           & --- \\
Matching subventions   & \code{grants.json}              & --- \\
Q\&A RAG               & \code{faq.json}                 & guides + régulations \\
\bottomrule
\end{tabularx}
\end{table}

Cette séparation nette entre données structurées et données longues justifie \textit{a posteriori} l'approche hybride~: les premières sont consommées par des modules déterministes, les secondes par le pipeline \ac{RAG} qui en extrait les passages pertinents à chaque requête.
